\documentclass[10pt,a4paper]{amsart}
\usepackage{amsmath,amssymb,amsthm,graphicx}
% Retained minimal preamble for the Stacks Project chapter more-morphisms.tex (More on Morphisms).
% Verbatim excerpt of the chapter preamble: Xy-pic setup (native preamble lines 14-18) and the native
% theorem declarations with the mathematical macros (native preamble lines 154-217). Copied mechanically;
% nothing added, removed or redefined. The wrapper class and the amsmath/amssymb/amsthm/graphicx base
% packages are supplied by the assembler, so the class/externaldocument part of the native preamble is omitted.
% For commutative diagrams we use Xy-pic
\usepackage[all]{xy}

% We use 2cell for 2-commutative diagrams.
\xyoption{2cell}
% Theorem environments.
%
\theoremstyle{plain}
\newtheorem{theorem}[subsection]{Theorem}
\newtheorem{proposition}[subsection]{Proposition}
\newtheorem{lemma}[subsection]{Lemma}

\theoremstyle{definition}
\newtheorem{definition}[subsection]{Definition}
\newtheorem{example}[subsection]{Example}
\newtheorem{exercise}[subsection]{Exercise}
\newtheorem{situation}[subsection]{Situation}

\theoremstyle{remark}
\newtheorem{remark}[subsection]{Remark}
\newtheorem{remarks}[subsection]{Remarks}

\numberwithin{equation}{subsection}

% Macros
%
\def\lim{\mathop{\mathrm{lim}}\nolimits}
\def\colim{\mathop{\mathrm{colim}}\nolimits}
\def\Spec{\mathop{\mathrm{Spec}}}
\def\Hom{\mathop{\mathrm{Hom}}\nolimits}
\def\Ext{\mathop{\mathrm{Ext}}\nolimits}
\def\SheafHom{\mathop{\mathcal{H}\!\mathit{om}}\nolimits}
\def\SheafExt{\mathop{\mathcal{E}\!\mathit{xt}}\nolimits}
\def\Sch{\mathit{Sch}}
\def\Mor{\mathop{\mathrm{Mor}}\nolimits}
\def\Ob{\mathop{\mathrm{Ob}}\nolimits}
\def\Sh{\mathop{\mathit{Sh}}\nolimits}
\def\NL{\mathop{N\!L}\nolimits}
\def\CH{\mathop{\mathrm{CH}}\nolimits}
\def\proetale{{pro\text{-}\acute{e}tale}}
\def\etale{{\acute{e}tale}}
\def\QCoh{\mathit{QCoh}}
\def\Ker{\mathop{\mathrm{Ker}}}
\def\Im{\mathop{\mathrm{Im}}}
\def\Coker{\mathop{\mathrm{Coker}}}
\def\Coim{\mathop{\mathrm{Coim}}}

% Boxtimes
%
\DeclareMathSymbol{\boxtimes}{\mathbin}{AMSa}{"02}

%
% Macros for moduli stacks/spaces
%
\def\QCohstack{\mathcal{QC}\!\mathit{oh}}
\def\Cohstack{\mathcal{C}\!\mathit{oh}}
\def\Spacesstack{\mathcal{S}\!\mathit{paces}}
\def\Quotfunctor{\mathrm{Quot}}
\def\Hilbfunctor{\mathrm{Hilb}}
\def\Curvesstack{\mathcal{C}\!\mathit{urves}}
\def\Polarizedstack{\mathcal{P}\!\mathit{olarized}}
\def\Complexesstack{\mathcal{C}\!\mathit{omplexes}}
% \Pic is the operator that assigns to X its picard group, usage \Pic(X)
% \Picardstack_{X/B} denotes the Picard stack of X over B
% \Picardfunctor_{X/B} denotes the Picard functor of X over B
\def\Pic{\mathop{\mathrm{Pic}}\nolimits}
\def\Picardstack{\mathcal{P}\!\mathit{ic}}
\def\Picardfunctor{\mathrm{Pic}}
\def\Deformationcategory{\mathcal{D}\!\mathit{ef}}
\newenvironment{reference}{}{}
\newenvironment{slogan}{}{}
\newenvironment{history}{}{}
\usepackage{cleveref}
\crefname{originallemma}{original lemma}{original lemmas}
\Crefname{originallemma}{Original lemma}{Original lemmas}
\makeatletter
\newlabel{descent-lemma-flat-surjective-quasi-compact-monomorphism-isomorphism}{{06NC}{source}{}{sourcepointer-1}{}}
\newlabel{descent-lemma-flat-surjective-quasi-compact-monomorphism-isomorphism@cref}{{[originallemma][0][]06NC}{[1][0][]source}{}{}{}}
\makeatother
\crefname{originallemma}{original lemma}{original lemmas}
\Crefname{originallemma}{Original lemma}{Original lemmas}
\makeatletter
\newlabel{morphisms-lemma-characterize-flat-closed-immersions}{{04PW}{source}{}{sourcepointer-2}{}}
\newlabel{morphisms-lemma-characterize-flat-closed-immersions@cref}{{[originallemma][0][]04PW}{[1][0][]source}{}{}{}}
\makeatother
\crefname{originallemma}{original lemma}{original lemmas}
\Crefname{originallemma}{Original lemma}{Original lemmas}
\makeatletter
\newlabel{morphisms-lemma-homeomorphism-affine}{{04DE}{source}{}{sourcepointer-3}{}}
\newlabel{morphisms-lemma-homeomorphism-affine@cref}{{[originallemma][0][]04DE}{[1][0][]source}{}{}{}}
\makeatother
\crefname{originallemma}{original lemma}{original lemmas}
\Crefname{originallemma}{Original lemma}{Original lemmas}
\makeatletter
\newlabel{morphisms-lemma-universally-injective}{{01S4}{source}{}{sourcepointer-4}{}}
\newlabel{morphisms-lemma-universally-injective@cref}{{[originallemma][0][]01S4}{[1][0][]source}{}{}{}}
\makeatother
\crefname{originallemma}{original lemma}{original lemmas}
\Crefname{originallemma}{Original lemma}{Original lemmas}
\makeatletter
\newlabel{schemes-lemma-monomorphism}{{01L3}{source}{}{sourcepointer-5}{}}
\newlabel{schemes-lemma-monomorphism@cref}{{[originallemma][0][]01L3}{[1][0][]source}{}{}{}}
\makeatother
\title{A weakly etale universal homeomorphism of schemes is an isomorphism}
\author{The Stacks Project Authors}
\date{stacks-project:more-morphisms.tex:local-snapshot-20261011; GNU FDL 1.2 or later}
\begin{document}
\maketitle
\noindent\textit{Editorial note.} Statement span is lines 19189-19193 of sources/stacks-more-morphisms/pinned-original-source.tex; proof span is lines 19195-19208, the complete author proof (12 body lines).\par\medskip
\noindent\textit{Editorial note.} Navigation stubs: 5 cited-but-not-included labels are supplied as non-load-bearing pointers with their pinned tags where the snapshot tag index has them.\par\medskip
\noindent\textit{Editorial note.} extra\_packages is empty on purpose: the retained preamble already loads xy with options and a second declaration clashes in the assembler wrapper.\par\medskip
\noindent\textit{Editorial note.} Text is retained verbatim, including the author's own typos and placeholders.\par\medskip
\noindent\textit{Editorial note.} Closing card of the weakly etale trilogy; together with L16089-L16090 it forms the chapter's self-contained weakly etale toolkit.\par\medskip
\noindent\textit{Allowed standard prerequisites (editorial).}\begin{itemize}
\item Explicit prerequisite boundary. Morphisms-chapter facts that a universal homeomorphism is submersive with universally injective underlying map and that a flat closed immersion is an open-and-closed immersion/isomorphism on components; Schemes, lemma monomorphism (diagonal iso iff monomorphism); Descent, lemma flat-surjective-quasi-compact-monomorphism-isomorphism.
\end{itemize}\medskip
\ifdefined\hypertarget\hypertarget{sourcepointer-1}{}\fi
\noindent\textit{Source reference.} Original lemma 06NC: sources/stacks-more-morphisms/pinned-original-source.tex, source lines 19208--19208. Cited by the retained statement or proof but not part of the collected spans; supplied only as a navigation pointer. The argument does not depend on its content.\par\medskip
\ifdefined\hypertarget\hypertarget{sourcepointer-2}{}\fi
\noindent\textit{Source reference.} Original lemma 04PW: sources/stacks-more-morphisms/pinned-original-source.tex, source lines 19208--19208. Cited by the retained statement or proof but not part of the collected spans; supplied only as a navigation pointer. The argument does not depend on its content.\par\medskip
\ifdefined\hypertarget\hypertarget{sourcepointer-3}{}\fi
\noindent\textit{Source reference.} Original lemma 04DE: sources/stacks-more-morphisms/pinned-original-source.tex, source lines 19208--19208. Cited by the retained statement or proof but not part of the collected spans; supplied only as a navigation pointer. The argument does not depend on its content.\par\medskip
\ifdefined\hypertarget\hypertarget{sourcepointer-4}{}\fi
\noindent\textit{Source reference.} Original lemma 01S4: sources/stacks-more-morphisms/pinned-original-source.tex, source lines 19208--19208. Cited by the retained statement or proof but not part of the collected spans; supplied only as a navigation pointer. The argument does not depend on its content.\par\medskip
\ifdefined\hypertarget\hypertarget{sourcepointer-5}{}\fi
\noindent\textit{Source reference.} Original lemma 01L3: sources/stacks-more-morphisms/pinned-original-source.tex, source lines 19208--19208. Cited by the retained statement or proof but not part of the collected spans; supplied only as a navigation pointer. The argument does not depend on its content.\par\medskip
\begin{lemma}
\label{lemma-weakly-etale-universal-homeomorphism}
Let $f : X \to Y$ be a morphism of schemes. If $f$ is weakly \'etale and
a universal homeomorphism, it is an isomorphism. 
\end{lemma}

\begin{proof}
Since $f$ is a universal homeomorphism, the diagonal
$\Delta : X \to X \times_Y X$ is a surjective closed immersion by
Morphisms, Lemmas \ref{morphisms-lemma-homeomorphism-affine} and
\ref{morphisms-lemma-universally-injective}. Since $\Delta$ is also
flat, we see that $\Delta$ must be an isomorphism by
Morphisms, Lemma \ref{morphisms-lemma-characterize-flat-closed-immersions}.
In other words, $f$ is a monomorphism
(Schemes, Lemma \ref{schemes-lemma-monomorphism}).
Since $f$ is a universal homeomorphism it is certainly
quasi-compact. Hence by Descent, Lemma
\ref{descent-lemma-flat-surjective-quasi-compact-monomorphism-isomorphism}
we find that $f$ is an isomorphism.
\end{proof}

\end{document}
